\documentclass{ximera}

\title{Elektrische krachtwerking - Inleiding}
\author{Daan Lipkens}

\begin{document}
\begin{abstract}
    Inleiding voor ladingen en kracht tussen ladingen.
\end{abstract}
\maketitle
% ----------------------------------------------------------------------
%          inleiding hoofdstuk
% ----------------------------------------------------------------------
\section{Inleiding}
De elektromagnetische kracht tussen geladen deeltjes is één van de vier fundamentele natuurkrachten. In dit hoofdstuk onderzoeken we eerst de basiseigenschappen van elektrische lading en de krachten die erop inwerken. Vervolgens bekijken we hoe objecten een lading kunnen krijgen, onder andere met behulp van de \emph{tribo-elektrische reeks}.

Daarna verdiepen we ons in de \emph{wet van Coulomb}, die zowel de grootte als de richting van de kracht tussen geladen deeltjes beschrijft. Je leert hoe je deze kracht kunt berekenen in uiteenlopende situaties, en we sluiten af met enkele praktische toepassingen.

Dit geheel behoort tot de \emph{elektrostatica}: het gedrag van stilstaande ladingen. In latere hoofdstukken bestuderen we ook bewegende ladingen en de verschijnselen die daarbij horen.
\begin{definition}
\emph{Elektrostatica} is het deel van de fysica dat het gedrag en de interacties van elektrische ladingen in rust en de krachtwerking er tussen bestudeert.
\end{definition}


\end{document}
